\documentclass[12pt, a4paper]{article}

% ---- Standard Packages ----
\usepackage[a4paper, margin=1in]{geometry}
\usepackage[utf8]{inputenc}
\usepackage[T1]{fontenc}
\usepackage[Indonesian]{babel}
\usepackage{amsmath, amssymb, amsthm}
\usepackage{enumitem}

\title{SOAL LATIHAN SEJARAH KELAS 10\\ \large Materi: Konsep Dasar, Metode Penelitian, Manusia Purba, dan Kerajaan Nusantara}
\author{Persiapan Ulangan Akhir Semester}
\date{}

\begin{document}
\maketitle

\newpage

\section*{Soal Pilihan Ganda}

\begin{enumerate}[label=\arabic*.]

    \item Kata \textit{sajaratun} dalam bahasa Arab yang menjadi asal-usul kata sejarah memiliki arti secara harfiah yaitu...
    \begin{enumerate}[label=\Alph*.]
        \item Ranting
        \item Akar
        \item Pohon
        \item Silsilah
        \item Daun
    \end{enumerate}

    \item Istilah \textit{history} dalam bahasa Inggris merujuk pada masa lampau umat manusia yang disusun berdasarkan...
    \begin{enumerate}[label=\Alph*.]
        \item Bukti tertulis maupun lisan
        \item Perkiraan sejarawan semata
        \item Cerita fiktif turun-temurun
        \item Mitos kuno masyarakat
        \item Hukum adat setempat
    \end{enumerate}

    \item Kata \textit{istoria} dari bahasa Yunani kuno memiliki makna dasar sebagai...
    \begin{enumerate}[label=\Alph*.]
        \item Cerita khayalan masa lalu
        \item Silsilah keturunan raja
        \item Hasil penelitian atau penyelidikan mendalam
        \item Hukum sebab-akibat alamiah
        \item Kitab suci keagamaan
    \end{enumerate}

    \item Tokoh yang dikenal sebagai Bapak Sejarah dan menyatakan bahwa sejarah bergerak layaknya garis lingkaran yang dipengaruhi oleh jatuh bangunnya peradaban adalah...
    \begin{enumerate}[label=\Alph*.]
        \item Aristoteles
        \item Ibnu Khaldun
        \item Sartono Kartodirdjo
        \item Herodotus
        \item E.H. Carr
    \end{enumerate}

    \item Menurut Aristoteles, sejarah merupakan sebuah sistem yang meneliti kejadian awal dan tersusun dalam bentuk...
    \begin{enumerate}[label=\Alph*.]
        \item Kronologi peristiwa dengan catatan rekaman bukti konkret serta empiris
        \item Cerita fiksi masa lampau
        \item Hukum alam yang pasti dan tidak dapat berubah
        \item Konstruksi imajinatif penulis sejarah
        \item Dongeng moral kemanusiaan
    \end{enumerate}

    \item Ibnu Khaldun mendefinisikan sejarah sebagai catatan mengenai manusia dan peradabannya yang mencakup...
    \begin{enumerate}[label=\Alph*.]
        \item Kronologi politik para penguasa
        \item Silsilah keluarga dinasti kerajaan secara berurutan
        \item Proses perubahan sosial, watak masyarakat, dan sebab-akibat secara filosofis
        \item Cerita mitologi tentang penciptaan bumi
        \item Rekonstruksi verbal berdasarkan ingatan kolektif
    \end{enumerate}

    \item Rekonstruksi masa lampau dalam bentuk cerita yang disusun kembali secara verbal berdasarkan fakta-fakta yang dapat dipertanggungjawabkan merupakan pengertian sejarah menurut...
    \begin{enumerate}[label=\Alph*.]
        \item Sartono Kartodirdjo
        \item Herodotus
        \item Aristoteles
        \item Ibnu Khaldun
        \item Kuntowijoyo
    \end{enumerate}

    \item Peristiwa atau kejadian nyata di masa lampau yang benar-benar terjadi tanpa adanya pengaruh dari opini, pandangan personal, atau prasangka peneliti dinamakan...
    \begin{enumerate}[label=\Alph*.]
        \item Sejarah subjektif
        \item Sejarah objektif
        \item Sejarah sebagai seni
        \item Tradisi lisan
        \item Dongeng masa lalu
    \end{enumerate}

    \item Sejarah sebagai peristiwa memiliki tiga karakteristik utama, yaitu...
    \begin{enumerate}[label=\Alph*.]
        \item Empiris, teoritis, dan memiliki metode
        \item Unik, abadi, dan penting
        \item Subjektif, imajinatif, dan emosional
        \item Lokal, nasional, dan global
        \item Tertulis, lisan, dan benda
    \end{enumerate}

    \item Penulisan sejarah yang membutuhkan keterlibatan intuisi, imajinasi, emosi, dan gaya bahasa yang indah agar pembaca merasakan atmosfer masa lalu termasuk dalam ruang lingkup...
    \begin{enumerate}[label=\Alph*.]
        \item Sejarah sebagai ilmu
        \item Sejarah sebagai seni
        \item Sejarah sebagai peristiwa
        \item Sejarah sebagai fakta keras
        \item Sejarah sebagai kronologi
    \end{enumerate}

    \item Kajian sejarah yang berfokus pada perkembangan kekuasaan, pemerintahan, kepemimpinan, konflik kekuasaan, dan diplomasi disebut...
    \begin{enumerate}[label=\Alph*.]
        \item Sejarah Perekonomian
        \item Sejarah Agama
        \item Sejarah Sosial
        \item Sejarah Politik
        \item Sejarah Kebudayaan
    \end{enumerate}

    \item Pembabakan waktu atau pembuatan babakan waktu dalam sejarah Indonesia dibuat bertujuan untuk...
    \begin{enumerate}[label=\Alph*.]
        \item Menentukan lokasi geografis peristiwa
        \item Memudahkan pemahaman terhadap babakan waktu sejarah
        \item Menghilangkan unsur subjektivitas sejarawan
        \item Mengumpulkan sumber sejarah primer
        \item Menyusun cerita fiktif masa lalu
    \end{enumerate}

    \item Sejarah selalu terikat oleh dua dimensi utama yang membatasi kapan dan di mana peristiwa berlangsung, yaitu...
    \begin{enumerate}[label=\Alph*.]
        \item Objek dan subjek
        \item Ruang (spasial) dan waktu (temporal)
        \item Fakta keras dan fakta lunak
        \item Lisan dan tulisan
        \item Teori dan metode
    \end{enumerate}

    \item Keterangan langsung dari pelaku atau saksi sejarah yang masih hidup mengenai peristiwa yang mereka alami sendiri disebut...
    \begin{enumerate}[label=\Alph*.]
        \item Tradisi lisan (\textit{Oral Tradition})
        \item Lisan (\textit{Oral History})
        \item Sumber Sekunder
        \item Artefak
        \item Fakta Mental
    \end{enumerate}

    \item Berikut ini yang \textbf{bukan} merupakan sifat dari sejarah sebagai ilmu adalah...
    \begin{enumerate}[label=\Alph*.]
        \item Empiris
        \item Memiliki Objek
        \item Memiliki Teori
        \item Memiliki Metode
        \item Bersifat Fiktif
    \end{enumerate}

    \item Menurut E.H. Carr, sejarah adalah suatu proses interaksi yang berkesinambungan dan sebuah dialog tanpa akhir antara...
    \begin{enumerate}[label=\Alph*.]
        \item Manusia dan alam sekitarnya
        \item Penguasa dan rakyat jelata
        \item Masa kini dan masa lalu
        \item Fakta keras dan fakta lunak
        \item Arkeologi dan paleontologi
    \end{enumerate}

    \item Elemen seni dalam penulisan sejarah di mana sejarawan menggunakan kemampuan membayangkan bagaimana suasana dan detail peristiwa masa lampau terjadi berdasarkan data disebut...
    \begin{enumerate}[label=\Alph*.]
        \item Intuisi
        \item Imajinasi
        \item Emosi
        \item Gaya Bahasa
        \item Objektivitas
    \end{enumerate}

    \item Langkah awal dalam metode penelitian sejarah di mana peneliti menentukan masalah atau tema yang akan diteliti dinamakan...
    \begin{enumerate}[label=\Alph*.]
        \item Heuristik
        \item Kritik Sumber
        \item Interpretasi
        \item Pemilihan Topik
        \item Historiografi
    \end{enumerate}

    \item Bekal utama peneliti berupa rasa suka, ketertarikan, dan motivasi yang memudahkan pengkajian objek penelitian jangka panjang disebut...
    \begin{enumerate}[label=\Alph*.]
        \item Kedekatan Intelektual
        \item Kedekatan Emosional
        \item Kritik Internal
        \item Sintesis Sejarah
        \item Intuisi Seni
    \end{enumerate}

    \item Tahapan pengumpulan data atau berburu sumber-sumber sejarah di lapangan, arsip, atau perpustakaan disebut...
    \begin{enumerate}[label=\Alph*.]
        \item Verifikasi
        \item Heuristik
        \item Interpretasi
        \item Historiografi
        \item Sintesis
    \end{enumerate}

    \item Proses menguji keaslian fisik dari sumber sejarah (seperti jenis kertas, tinta, atau bahan prasasti) dinamakan...
    \begin{enumerate}[label=\Alph*.]
        \item Kritik Eksternal
        \item Kritik Internal
        \item Analisis
        \item Sintesis
        \item Historiografi
    \end{enumerate}

    \item Proses melakukan pengujian terhadap tingkat kepercayaan atau kredibilitas dari isi informasi suatu sumber sejarah disebut...
    \begin{enumerate}[label=\Alph*.]
        \item Kritik Eksternal
        \item Kritik Internal
        \item Heuristik
        \item Pemilihan Topik
        \item Sintesis
    \end{enumerate}

    \item Dalam tahap interpretasi, proses mengurai fakta sejarah untuk memperoleh penjelasan dari sumber yang tidak secara gamblang menjelaskan kejadian disebut...
    \begin{enumerate}[label=\Alph*.]
        \item Analisis
        \item Sintesis
        \item Verifikasi
        \item Heuristik
        \item Narasi
    \end{enumerate}

    \item Proses menyatukan, menghubungkan, dan merangkai hasil analisis-analisis dari berbagai sumber sejarah guna mencapai kesimpulan yang utuh dinamakan...
    \begin{enumerate}[label=\Alph*.]
        \item Analisis
        \item Heuristik
        \item Kritik Sumber
        \item Sintesis
        \item Historiografi
    \end{enumerate}

    \item Tahapan terakhir dalam metode sejarah yang berupa proses penulisan fakta-fakta menjadi karya ilmiah yang utuh dan sistematis dinamakan...
    \begin{enumerate}[label=\Alph*.]
        \item Interpretasi
        \item Verifikasi
        \item Heuristik
        \item Historiografi
        \item Kronologi
    \end{enumerate}

    \item Pendekatan khas ilmu sejarah yang memiliki karakteristik memanjang dalam waktu tetapi menyempit dalam ruang dinamakan...
    \begin{enumerate}[label=\Alph*.]
        \item Pendekatan Sinkronis
        \item Pendekatan Diakronis
        \item Pendekatan Spasial
        \item Pendekatan Struktural
        \item Pendekatan Statis
    \end{enumerate}

    \item Pendekatan penelitian sejarah yang meluas dalam ruang tetapi terbatas/menyempit dalam waktu serta bersifat struktural dinamakan...
    \begin{enumerate}[label=\Alph*.]
        \item Pendekatan Diakronis
        \item Pendekatan Sinkronis
        \item Pendekatan Kronologis
        \item Pendekatan Temporal
        \item Pendekatan Empiris
    \end{enumerate}

    \item Ilmu bantu sejarah yang mengkaji bekas warisan masa lalu berupa benda fisik purba atau perkakas buatan manusia (artefak) adalah...
    \begin{enumerate}[label=\Alph*.]
        \item Paleontologi
        \item Arkeologi
        \item Epigrafi
        \item Filologi
        \item Numismatik
    \end{enumerate}

    \item Disiplin ilmu bantu yang mengkaji bentuk kehidupan purba, perkembangan flora, dan fauna kuno berdasarkan sisa-sisa organik berupa fosil disebut...
    \begin{enumerate}[label=\Alph*.]
        \item Paleoantropologi
        \item Paleografi
        \item Paleontologi
        \item Keramologi
        \item Ikonografi
    \end{enumerate}

    \item Cabang ilmu bantu sejarah yang mempelajari tulisan kuno yang dipahat pada material keras seperti batu (prasasti) atau logam kuno dinamakan...
    \begin{enumerate}[label=\Alph*.]
        \item Paleografi
        \item Filologi
        \item Epigrafi
        \item Ikonografi
        \item Numismatik
    \end{enumerate}

    \item Ilmu bantu sejarah yang mengkaji mata uang kuno, medalion, atau segel perdagangan masa lalu dinamakan...
    \begin{enumerate}[label=\Alph*.]
        \item Numismatik
        \item Keramologi
        \item Genealogi
        \item Epigrafi
        \item Paleografi
    \end{enumerate}

    \item Ilmu bantu filologi sangat berguna bagi sejarawan untuk melakukan analisis yang berkaitan dengan...
    \begin{enumerate}[label=\Alph*.]
        \item Struktur geologi tempat fosil berada
        \item Asal-usul keramik dan porselen peninggalan asing
        \item Bahasa, teks sastra kuno, naskah manuskrip lama, dan perubahan kebahasaan
        \item Silsilah silsilah keturunan raja
        \item Simbol, arca, dan relief keagamaan kuno
    \end{enumerate}

    \item Ketika manusia berperan aktif merekonstruksi dan menulis kembali sebuah kejadian masa lalu, maka dalam dimensi sejarah manusia berposisi sebagai...
    \begin{enumerate}[label=\Alph*.]
        \item Objek sejarah
        \item Subjek sejarah
        \item Saksi pasif
        \item Ilmu bantu sejarah
        \item Fakta keras
    \end{enumerate}

    \item Menurut sejarawan Kuntowijoyo, konsep waktu dalam sejarah mencakup 4 dimensi utama, yaitu...
    \begin{enumerate}[label=\Alph*.]
        \item Masa lalu, masa kini, masa depan, dan masa prasejarah
        \item Perkembangan, kesinambungan, pengulangan, dan perubahan
        \item Manusia, ruang, waktu, dan lingkungan
        \item Heuristik, verifikasi, interpretasi, dan historiografi
        \item Politik, ekonomi, sosial, dan budaya
    \end{enumerate}

    \item Dimensi waktu yang terjadi jika masyarakat baru tetap mengadopsi atau melanjutkan lembaga, kebiasaan, atau tradisi lama dinamakan...
    \begin{enumerate}[label=\Alph*.]
        \item Perkembangan
        \item Kesinambungan
        \item Pengulangan
        \item Perubahan
        \item Keberlanjutan Total
    \end{enumerate}

    \item Konsep pengulangan dalam dimensi waktu sejarah terjadi apabila...
    \begin{enumerate}[label=\Alph*.]
        \item Tokoh yang sama melakukan aktivitas yang sama di dua tempat berbeda
        \item Peristiwa di masa lalu kembali terjadi di masa sesudahnya dengan pola atau karakteristik mirip, meski pelaku dan waktu berbeda
        \item Masyarakat mengalami pergeseran secara besar-besaran dalam waktu singkat
        \item Perkembangan masyarakat berjalan stagnan tanpa adanya kemajuan sedikit pun
        \item Sejarawan menulis buku sejarah yang sama secara berulang kali
    \end{enumerate}

    \item Sumber sejarah yang diperoleh dari benda peninggalan yang memuat aksara atau tulisan kuno disebut...
    \begin{enumerate}[label=\Alph*.]
        \item Sumber Lisan
        \item Sumber Tertulis
        \item Sumber Audio-Visual
        \item Fakta Sosial
        \item Fakta Mental
    \end{enumerate}

    \item Prasasti, arsip resmi negara, dan kesaksian langsung dari pelaku sejarah yang berada di lokasi kejadian termasuk dalam kategori...
    \begin{enumerate}[label=\Alph*.]
        \item Sumber Primer
        \item Sumber Sekunder
        \item Sumber Tersier
        \item Fakta Lunak
        \item Tradisi Lisan
    \end{enumerate}

    \item Buku teks sejarah, biografi tokoh, dan ensiklopedia merupakan contoh konkret dari peninggalan berupa...
    \begin{enumerate}[label=\Alph*.]
        \item Sumber Primer
        \item Sumber Sekunder
        \item Artefak Fisik
        \item Fakta Mental
        \item Mentofak
    \end{enumerate}

    \item Kesimpulan akhir atau hasil analisis kritis yang diperoleh dari pengolahan berbagai bukti-bukti sejarah disebut...
    \begin{enumerate}[label=\Alph*.]
        \item Bukti Sejarah
        \item Fakta Sejarah
        \item Sumber Primer
        \item Artefak
        \item Historiografi
    \end{enumerate}

    \item Kondisi psikologis, keyakinan, ideologi, alam pikiran, dan tatanan nilai mental yang berkembang dalam masyarakat periode tertentu disebut...
    \begin{enumerate}[label=\Alph*.]
        \item Artefak
        \item Fakta Sosial (Sosiofak)
        \item Fakta Mental (Mentofak)
        \item Fakta Keras
        \item Fakta Lunak
    \end{enumerate}

    \item Fakta sejarah yang kebenarannya sudah mutlak, pasti, teruji, dan tidak dapat diperdebatkan lagi oleh siapa pun (contoh: Proklamasi 17 Agustus 1945) dinamakan...
    \begin{enumerate}[label=\Alph*.]
        \item Fakta Lunak (\textit{Soft Fact})
        \item Fakta Keras (\textit{Hard Fact})
        \item Fakta Opini
        \item Fakta Mental
        \item Fakta Subjektif
    \end{enumerate}

    \item Pernyataan mengenai kepastian letak pusat Kerajaan Sriwijaya yang hingga kini masih terbuka untuk diteliti karena minimnya bukti pendukung merupakan contoh...
    \begin{enumerate}[label=\Alph*.]
        \item Fakta Keras (\textit{Hard Fact})
        \item Fakta Lunak (\textit{Soft Fact})
        \item Artefak Fisik
        \item Sosiofak
        \item Mentofak
    \end{enumerate}

    \item Manusia purba \textit{Meganthropus paleojavanicus} ditemukan oleh G.H.R. von Koenigswald pada tahun 1936 dan 1941 di daerah...
    \begin{enumerate}[label=\Alph*.]
        \item Trinil
        \item Sangiran
        \item Mojokerto
        \item Ngandong
        \item Liang Bua
    \end{enumerate}

    \item Manusia purba jenis \textit{Pithecanthropus} yang memiliki usia paling tua dan ditemukan di Kepuhklagen oleh Von Koenigswald adalah...
    \begin{enumerate}[label=\Alph*.]
        \item \textit{Pithecanthropus erectus}
        \item \textit{Pithecanthropus soloensis}
        \item \textit{Pithecanthropus mojokertensis}
        \item \textit{Homo wajakensis}
        \item \textit{Homo floresiensis}
    \end{enumerate}

    \item Fosil manusia purba \textit{Pithecanthropus erectus} yang berarti "manusia kera yang berjalan tegak" ditemukan oleh Eugene Dubois di daerah...
    \begin{enumerate}[label=\Alph*.]
        \item Sangiran
        \item Ngandong
        \item Sambungmacan
        \item Trinil, Ngawi
        \item Wajak, Tulungagung
    \end{enumerate}

    \item Manusia purba dari Pulau Flores yang ditemukan tahun 2003 dengan karakteristik fisik kerdil (tinggi sekitar 1 meter) dan volume otak 380 cc bernama...
    \begin{enumerate}[label=\Alph*.]
        \item \textit{Homo soloensis}
        \item \textit{Homo sapiens}
        \item \textit{Homo wajakensis}
        \item \textit{Homo floresiensis}
        \item \textit{Pithecanthropus erectus}
    \end{enumerate}

    \item Pembabakan zaman praaksara secara arkeologis menurut Prof. Dr. R. Soekmono dibagi menjadi dua periodisasi besar, yaitu...
    \begin{enumerate}[label=\Alph*.]
        \item Zaman Berburu dan Zaman Bercocok Tanam
        \item Zaman Batu dan Zaman Logam
        \item Zaman Primer dan Zaman Sekunder
        \item Zaman Undagi dan Zaman Melayu
        \item Zaman Praaksara dan Zaman Aksara
    \end{enumerate}

    \item Perubahan pola hidup masyarakat praaksara dari mencari makanan (\textit{food gathering}) menjadi memproduksi makanan sendiri (\textit{food producing}) terjadi pada fase revolusi kebudayaan zaman...
    \begin{enumerate}[label=\Alph*.]
        \item Paleolitikum
        \item Mesolitikum
        \item Neolitikum
        \item Megalitikum
        \item Zaman Besi
    \end{enumerate}

    \item Pola hidup menetap secara permanen yang mulai diadopsi oleh masyarakat praaksara pada masa Neolitikum dikenal dengan istilah...
    \begin{enumerate}[label=\Alph*.]
        \item Nomaden
        \item Sedenter
        \item Semisedenter
        \item Migrasi
        \item Kolonisasi
    \end{enumerate}

    \item Munculnya golongan "Undagi" pada masa akhir praaksara menandakan adanya masa perkampungan teratur dengan keahlian khusus yaitu...
    \begin{enumerate}[label=\Alph*.]
        \item Membuat rakit untuk pelayaran
        \item Menjinakkan hewan liar menjadi ternak
        \item Pembuatan peralatan dan perhiasan dari bahan logam
        \item Berburu binatang menggunakan batu kasar
        \item Menulis prasasti pada material batu keras
    \end{enumerate}

    \item Teknik melebur dan mencetak bijih logam dengan menggunakan cetak batu dinamakan teknik...
    \begin{enumerate}[label=\Alph*.]
        \item \textit{A cire perdue}
        \item \textit{Bivalve}
        \item \textit{Flakes}
        \item \textit{Chopper}
        \item \textit{Pebble}
    \end{enumerate}

    \item Suku Dayak, Toraja, Nias, dan Batak di Indonesia merupakan keturunan langsung dari gelombang migrasi ras...
    \begin{enumerate}[label=\Alph*.]
        \item Proto Melayu
        \item Deutro Melayu
        \item Melanesoid
        \item Weddid
        \item Negroid
    \end{enumerate}

    \item Rute migrasi ras Deutro Melayu (Melayu Muda) ke Nusantara membawa kebudayaan logam yang berasal dari Vietnam, yaitu kebudayaan...
    \begin{enumerate}[label=\Alph*.]
        \item Kebudayaan Pacitan
        \item Kebudayaan Ngandong
        \item Kebudayaan Toala
        \item Kebudayaan Sampung
        \item Kebudayaan Dongson
    \end{enumerate}

    \item Sumber sejarah utama Kerajaan Kutai Martadipura yang berbentuk 7 buah tugu batu pemujaan ditulis menggunakan huruf Pallawa dan bahasa Sanskerta bernama...
    \begin{enumerate}[label=\Alph*.]
        \item Prasasti Tugu
        \item Prasasti Ciaruteun
        \item \textit{Yupa}
        \item Kitab Sutasoma
        \item Kitab Negarakertagama
    \end{enumerate}

    \item Peninggalan budaya Pacitan pada zaman Paleolitikum yang berupa kapak genggam kasar dalam dunia arkeologi dikenal sebagai tradisi...
    \begin{enumerate}[label=\Alph*.]
        \item Tradisi serpih bilah (\textit{Flakes})
        \item Tradisi kapak perimbas-penetak (\textit{Chopper-chopping tool})
        \item Tradisi alat tulang dan tanduk
        \item Tradisi batu besar (\textit{Megalitikum})
        \item Tradisi cetak lilin (\textit{A cire perdue})
    \end{enumerate}

    \item Prasasti peninggalan Kerajaan Tarumanegara yang berisi informasi mengenai penggalian Sungai Candrabhaga dan Gomati sepanjang 11 km untuk irigasi dan mencegah banjir adalah...
    \begin{enumerate}[label=\Alph*.]
        \item Prasasti Ciaruteun
        \item Prasasti Jambu
        \item Prasasti Kebon Kopi
        \item Prasasti Tugu
        \item Prasasti Telaga Batu
    \end{enumerate}

    \item Kerajaan maritim bernuansa Buddha terbesar di Asia Tenggara yang menguasai jalur perdagangan internasional vital di Selat Malaka pada abad ke-7 hingga ke-14 Masehi adalah...
    \begin{enumerate}[label=\Alph*.]
        \item Kerajaan Kutai
        \item Kerajaan Tarumanegara
        \item Kerajaan Sriwijaya
        \item Kerajaan Majapahit
        \item Kesultanan Demak
    \end{enumerate}

    \item Ikrar politik ikonik bernama Sumpah Palapa yang bertujuan menyatukan seluruh wilayah Nusantara di bawah naungan Kerajaan Majapahit diucapkan oleh...
    \begin{enumerate}[label=\Alph*.]
        \item Raden Wijaya
        \item Hayam Wuruk
        \item Gajah Mada
        \item Kertarajasa
        \item Empu Prapanca
    \end{enumerate}

    \item Kerajaan Islam pertama di Nusantara yang berdiri sekitar abad ke-13 Masehi (1267 M) di pesisir utara Pulau Sumatra adalah...
    \begin{enumerate}[label=\Alph*.]
        \item Kesultanan Samudera Pasai
        \item Kesultanan Demak
        \item Kesultanan Mataram Islam
        \item Kesultanan Gowa-Tallo
        \item Kesultanan Banten
    \end{enumerate}
\item Kerajaan Kutai Martadipura yang merupakan kerajaan Hindu tertua di Nusantara secara geografis terletak di kawasan...
    \begin{enumerate}[label=\Alph*.]
        \item Muara Cianten, di tepi hulu Sungai Cisadane
        \item Muara Kaman, di tepi hulu Sungai Mahakam
        \item Trowulan, di dekat aliran Sungai Brantas
        \item Lhokseumawe, di pesisir hulu Sungai Pasai
        \item Kotagede, di dekat aliran Sungai Opak
    \end{enumerate}

    \item Berdasarkan catatan pada prasasti Yupa, Raja Mulawarman menyedekahkan 20.000 ekor sapi kepada kaum Brahmana di tanah suci pemujaan Dewa Siwa yang bernama...
    \begin{enumerate}[label=\Alph*.]
        \item Sriksetra
        \item Candrabhaga
        \item Waprakeswara
        \item Trowulan
        \item Alas Mentaok
    \end{enumerate}

    \item Kerajaan Tarumanegara merupakan kerajaan bercorak Hindu tertua di Pulau Jawa yang pusat kekuasaannya mencakup wilayah...
    \begin{enumerate}[label=\Alph*.]
        \item Jawa Timur, sekitar Mojokerto dan Kediri
        \item Jawa Tengah, sekitar Surakarta dan Yogyakarta
        \item Kalimantan Timur, sekitar Sungai Mahakam
        \item Jawa Barat, sekitar Bogor, Bekasi, Karawang, dan Jakarta
        \item Sumatra Selatan, sekitar Palembang
    \end{enumerate}

    \item Prasasti Tugu mengisahkan kepedulian ekonomi Raja Purnawarman terhadap rakyatnya melalui proyek penggalian dua saluran air sepanjang 11 km untuk irigasi dan mencegah banjir, yaitu Sungai...
    \begin{enumerate}[label=\Alph*.]
        \item Mahakam dan Barito
        \item Candrabhaga dan Gomati
        \item Bengawan Solo dan Brantas
        \item Ciaruteun dan Cidanghiang
        \item Pasai dan Malaka
    \end{enumerate}

    \item Keberadaan Kerajaan Sriwijaya sebagai pusat pembelajaran agama Buddha yang besar di luar India dikonfirmasi oleh catatan musafir Cina yang bernama...
    \begin{enumerate}[label=\Alph*.]
        \item Fa-Hien
        \item Ma Huan
        \item I-Tsing
        \item Cheng Ho
        \item Tome Pires
    \end{enumerate}

    \item Selain menguasai jalur perdagangan maritim, Kerajaan Sriwijaya juga dikenal luas di dunia internasional bertindak sebagai pusat kajian ilmiah agama Buddha aliran...
    \begin{enumerate}[label=\Alph*.]
        \item Hinayana
        \item Mahayana murni
        \item Vajrayana
        \item Waisnawa
        \item Tantrayana
    \end{enumerate}

    \item Pada abad ke-9 M, Kerajaan Sriwijaya mencapai puncak kejayaan politik dan ekonomi di bawah pemerintahan Raja Balaputradewa yang menjalin hubungan internasional erat dengan kerajaan India kuno, yaitu...
    \begin{enumerate}[label=\Alph*.]
        \item Kerajaan Chola
        \item Kerajaan Nalanda
        \item Kerajaan Majapahit
        \item Dinasti Tang
        \item Dinasti Yuan
    \end{enumerate}

    \item Raden Wijaya berhasil mendirikan Kerajaan Majapahit pada tahun 1293 M setelah memanfaatkan kedatangan pasukan asing untuk menghancurkan Jayakatwang dari Kediri. Pasukan asing tersebut berasal dari...
    \begin{enumerate}[label=\Alph*.]
        \item Pasukan Portugis Barat
        \item Pasukan Islam Samudera Pasai
        \item Pasukan Mongol (Dinasti Yuan)
        \item Pasukan Maritim Sriwijaya
        \item Pasukan Kesultanan Demak
    \end{enumerate}

    \item Saat mencapai puncak kejayaannya pada abad ke-14 M di bawah pemerintahan Raja Hayam Wuruk, sang raja memimpin dengan gelar resmi yaitu...
    \begin{enumerate}[label=\Alph*.]
        \item Kertarajasa Jayawardhana
        \item Sri Rajasanagara
        \item Sultan Malik As-Saleh
        \item Panembahan Senopati
        \item Sultan Trenggono
    \end{enumerate}

    \item Struktur wilayah dan sistem pemerintahan Kerajaan Majapahit secara rinci termuat di dalam salah satu sumber sejarah tertulis, yaitu...
    \begin{enumerate}[label=\Alph*.]
        \item Kitab \textit{Sutasoma} karya Empu Tantular
        \item Kitab \textit{Suma Oriental} karya Tome Pires
        \item Kitab \textit{Negarakertagama} karya Empu Prapanca
        \item Kitab \textit{Babad Tanah Jawi}
        \item Prasasti Telaga Batu
    \end{enumerate}

    \item Pendiri Kesultanan Samudera Pasai yang naik takhta setelah memeluk agama Islam berkat bimbingan Syekh Ismail dari Makkah memiliki nama asli Nusantara yaitu...
    \begin{enumerate}[label=\Alph*.]
        \item Marah Silu
        \item Kudungga
        \item Raden Patah
        \item Raden Mas Sutawijaya
        \item Daeng Manrabia
    \end{enumerate}

    \item Sebagai bandar transito komersial internasional yang penting di Selat Malaka, Samudera Pasai menjadi kerajaan pertama di Nusantara yang mencetak mata uang emas sendiri yang disebut...
    \begin{enumerate}[label=\Alph*.]
        \item Dirham
        \item Dinor
        \item Kupang
        \item Real
        \item Saka
    \end{enumerate}

    \item Penjelajah terkenal asal Maroko yang mencatat bahwa Sultan Samudera Pasai adalah sosok yang sangat saleh, rendah hati, dan suka berdiskusi ilmiah mengenai hukum Islam dengan ulama adalah...
    \begin{enumerate}[label=\Alph*.]
        \item Marcopolo
        \item Tome Pires
        \item Ma Huan
        \item Ibnu Battutah
        \item I-Tsing
    \end{enumerate}

    \item Kesultanan Demak didirikan pada akhir abad ke-15 M oleh Raden Patah, yang secara silsilah keturunan masih merupakan putra dari raja terakhir Majapahit, yaitu...
    \begin{enumerate}[label=\Alph*.]
        \item Jayakatwang
        \item Raden Wijaya
        \item Hayam Wuruk
        \item Brawijaya V
        \item Sultan Trenggono
    \end{enumerate}

    \item Nama "Jayakarta" yang digunakan untuk menggantikan nama wilayah Sunda Kelapa setelah direbut dari pengaruh Portugis pada tahun 1527 diberikan oleh panglima perang Demak bernama...
    \begin{enumerate}[label=\Alph*.]
        \item Raden Patah
        \item Fatahillah
        \item Sunan Kalijaga
        \item Sultan Trenggono
        \item Gajah Mada
    \end{enumerate}

    \item Berdirinya Kesultanan Mataram Islam pada tahun 1588 M diawali dari sebuah peristiwa pembukaan wilayah hutan hadiah dari Kesultanan Pajang yang bernama...
    \begin{enumerate}[label=\Alph*.]
        \item Alas Mentaok
        \item Taman Sriksetra
        \item Kawasan Muara Kaman
        \item Situs Trowulan
        \item Lembah Bengawan Solo
    \end{enumerate}

    \item Pembuatan Kalender Jawa, yang merupakan perpaduan harmonis antara kalender Saka dan kalender Hijriah, berhasil diciptakan pada masa kejayaan Mataram Islam oleh...
    \begin{enumerate}[label=\Alph*.]
        \item Ki Ageng Pemanahan
        \item Panembahan Senopati
        \item Sultan Agung Hanyokrokusumo
        \item Sultan Trenggono
        \item Raden Mas Sutawijaya
    \end{enumerate}

    \item Sebagai bentuk perlawanan terhadap kolonialisme Barat di Nusantara, Sultan Agung memimpin serangan militer kolosal Mataram Islam sebanyak dua kali ke Batavia pada tahun...
    \begin{enumerate}[label=\Alph*.]
        \item 1293 dan 1300
        \item 1527 dan 1530
        \item 1628 dan 1629
        \item 1667 dan 1669
        \item 1755 dan 1760
    \end{enumerate}

    \item Karakteristik unik dari sistem pemerintahan Kesultanan Gowa-Tallo di Sulawesi Selatan yang membedakannya dari kesultanan lain di Nusantara adalah...
    \begin{enumerate}[label=\Alph*.]
        \item Dipimpin langsung oleh kaum Brahmana dari India kuno
        \item Roda pemerintahan dikelola bersama hasil penggabungan dua kerajaan kembar
        \item Sistem kekuasaan mutlak berada di bawah koordinasi Wali Songo
        \item Tidak memiliki wilayah laut dan berfokus penuh pada agraris pedalaman
        \item Struktur kepemimpinannya dipimpin oleh musafir dari Tiongkok
    \end{enumerate}

    \item Dominasi maritim Kesultanan Gowa-Tallo sebagai bandar niaga internasional bebas (*free port*) akhirnya patah setelah VOC Belanda memaksa Sultan Hasanuddin menandatangani...
    \begin{enumerate}[label=\Alph*.]
        \item Perjanjian Giyanti (1755)
        \item Perjanjian Bongaya (1667)
        \item Sumpah Palapa (Abad ke-14)
        \item Piagam Trowulan (1293)
        \item Kapitulasi Tuntang (1811)
    \end{enumerate}
\end{enumerate}

\end{document}
